\chapter{Conclusiones y Trabajos Futuros}
\label{cap:conclusiones_trabajos_futuros}

En este capítulo final se exponen las conclusiones derivadas del diseño, desarrollo e implementación de la plataforma \textit{SkiClubManager}, evaluando el cumplimiento de los objetivos planteados inicialmente. Asimismo, se definen de forma integrada las líneas de investigación y desarrollo futuro recomendadas para la evolución del sistema en el contexto de la gestión técnica de clubes de esquí alpino.

\section{Conclusiones y líneas de desarrollo futuro}

El desarrollo de la plataforma \textit{SkiClubManager} ha permitido integrar de forma exitosa las tareas operativas críticas de un club de esquí en una única solución tecnológica centralizada. Al unificar la planificación de entrenamientos, la gestión logística del transporte, el control de asistencia y el almacenamiento multimedia, el sistema mitiga la dependencia estructural de herramientas analógicas e informales dispersas. Desde una perspectiva técnico-metodológica, la combinación de una \textit{API REST} en \textit{Node.js} y un motor relacional \textit{MySQL} ha demostrado ser una arquitectura sólida y mantenible, capaz de garantizar la consistencia de los datos y gestionar operaciones concurrentes críticas, especialmente en los registros de asistencia y asignación de plazas de transporte para los deportistas.

Por su parte, la capa de interacción móvil implementada con \textit{React Native} y soportada por la \textit{Context API} valida la viabilidad del desarrollo multiplataforma bajo una única base de código, optimizando los tiempos de cómputo y ofreciendo una interfaz fluida y adaptada a las condiciones climáticas de alta montaña. Asimismo, la seguridad del ecosistema queda blindada mediante el estándar de autenticación basado en \textit{JSON Web Tokens} (\textit{JWT}), el cual restringe el acceso a los recursos según los roles validados del \textit{backend}. Finalmente, la integración de la librería de analítica visual transforma los registros históricos planos en gráficos estadísticos comprensibles, aportando un valor pedagógico y operativo directo tanto para el cuerpo técnico como para las familias de los deportistas.

A pesar de que la aplicación cumple rigurosamente con el alcance funcional establecido para este proyecto de ingeniería, la flexibilidad de la arquitectura \textit{online} actual —diseñada de manera desacoplada mediante controladores y rutas independientes— constituye una base tecnológica ya preparada para absorber nuevas extensiones del sistema sin alterar el núcleo del negocio. Aprovechando esta infraestructura distribuida, una de las líneas prioritarias de trabajo futuro es la implementación de un módulo para que los entrenadores elaboren informes cualitativos de evaluación técnica en pista. Al estar el servicio \textit{backend} completamente operativo en la nube, la incorporación de este módulo requerirá únicamente exponer nuevos \textit{endpoints} relacionales para estructurar dichas evaluaciones, facilitando que el cuerpo técnico envíe los reportes de rendimiento en tiempo real a la base de datos centralizada.

Esta capacidad de expansión en línea se complementará, en el plano funcional, con la inclusión del rol nativo del deportista, permitiendo al esquiador interactuar de forma directa con su propio histórico, entrenamientos y progresión técnica. Asimismo, debido a la naturaleza extrema del entorno geográfico alpino, se proyecta evolucionar la aplicación móvil hacia una arquitectura híbrida \textit{Offline-First} mediante motores de persistencia local como \textit{SQLite}, facultando al equipo para trabajar en zonas de pista con conectividad nula y sincronizar los informes una vez restablecido el acceso a la red. Por último, se contempla la ampliación de las capacidades de almacenamiento para soportar archivos de vídeo pesados orientados al análisis técnico, la integración de notificaciones \textit{push} a través de \textit{Firebase Cloud Messaging} para alertas y la implementación de una pasarela de pago como \textit{Stripe} para centralizar el recaudo de cuotas institucionales.